% =========================================================================
% Impact of Hybrid Structured Random Projections on GNN Performance
% Final Research Paper — IEEE Conference Format
% University of Peradeniya, Dept. of Computer Engineering
% =========================================================================
\documentclass[conference]{IEEEtran}

% --- Packages ---
\usepackage{cite}
\usepackage{amsmath,amssymb,amsfonts}
\usepackage{algorithmic}
\usepackage{algorithm}
\usepackage{graphicx}
\usepackage{textcomp}
\usepackage{xcolor}
\usepackage{booktabs}
\usepackage{multirow}
\usepackage[hidelinks]{hyperref}

% --- IEEEtran author line-break helper ---
\makeatletter
\newcommand{\linebreakand}{%
  \end{@IEEEauthorhalign}
  \hfill\mbox{}\par
  \mbox{}\hfill\begin{@IEEEauthorhalign}
}
\makeatother

% =========================================================================
\begin{document}

\title{Impact of Hybrid Structured Random Projections on Graph Embeddings}

\author{
    \IEEEauthorblockN{H.A.D.T.N. Haththella}
    \IEEEauthorblockA{\textit{Dept. of Computer Engineering} \\
    \textit{University of Peradeniya}\\
    Sri Lanka \\
    e20133@eng.pdn.ac.lk}
    \and
    \IEEEauthorblockN{L.R.S.R. Premasiri}
    \IEEEauthorblockA{\textit{Dept. of Computer Engineering} \\
    \textit{University of Peradeniya}\\
    Sri Lanka \\
    e20305@eng.pdn.ac.lk}
    \and
    \IEEEauthorblockN{A.N.T. Somathilaka}
    \IEEEauthorblockA{\textit{Dept. of Computer Engineering} \\
    \textit{University of Peradeniya}\\
    Sri Lanka \\
    e20381@eng.pdn.ac.lk}
    \linebreakand
    \IEEEauthorblockN{Sampath Deegalla}
    \IEEEauthorblockA{\textit{Dept. of Computer Engineering} \\
    \textit{University of Peradeniya}\\
    Sri Lanka \\
    sampath@eng.pdn.ac.lk}
}

\maketitle

% =========================================================================
\begin{abstract}
Graph representation learning maps nodes to low-dimensional vectors for downstream tasks. Many scalable projection-based approaches employ random matrices to avoid the computational overhead of iterative message passing. While the standard FastRP framework achieves scalability through very sparse projections, alternative projection designs offer varying trade-offs between efficiency and embedding quality. This work investigates whether structured and hybrid random projections can outperform a dense Gaussian baseline within a FastRP-style iterative propagation pipeline.

For node classification, three projection strategies are compared: (i) a standard dense Gaussian projection as the baseline, (ii) a striped structured sparse projection, and (iii) a hybrid projection combining structural and node feature components. Link prediction experiments specifically contrast the Gaussian baseline against the hybrid strategy. Evaluations span seven benchmark datasets with Optuna-based hyperparameter tuning.

Results demonstrate that integrating node features into the initial projection substantially improves embedding quality on feature-rich networks. On the large-scale ogbn-arxiv benchmark, the hybrid approach yields a 16.8 percentage point accuracy gain over the Gaussian baseline. The striped structured projection reduces matrix density and remains competitive on sparse citation graphs, but exhibits sensitivity to graph topology on denser datasets such as Flickr. Overall, carefully designed structured and hybrid projections provide a compelling, dataset-aware alternative to standard random projections for scalable graph embedding.
\end{abstract}

\begin{IEEEkeywords}
Graph Embedding, FastRP, Random Projection, Hybrid Projection, Node
Classification, Link Prediction, Graph Neural Networks, Optuna
\end{IEEEkeywords}

% =========================================================================
\section{Introduction}
\label{sec:intro}

The growing prevalence of graph-structured data in social networks,
knowledge graphs, and citation databases has driven significant interest in
graph representation learning. The central goal is to map each node to a
low-dimensional vector that captures structural and semantic properties of
the graph, enabling effective downstream tasks such as node classification
and link prediction \cite{khemani2024, gkarmpounis2024}.

Among the available approaches, Graph Neural Networks (GNNs) have become a
dominant paradigm for learning node representations through iterative message
passing over local neighborhoods \cite{khemani2024}. However, end-to-end
trained GNNs face scalability challenges on large real-world graphs: full
neighborhood aggregation incurs memory and compute costs that scale with
graph size, and gradient-based training across many epochs amplifies these
costs \cite{bui2023rap}.

Projection-based graph embedding methods offer a compelling alternative. By
initializing node representations with random projections and propagating
them across the graph, methods such as FastRP \cite{chen2019fast}, RandNE
\cite{zhang2018randne}, and SketchNE \cite{xie2023sketchne} achieve
remarkable scalability without incurring the training overhead of deep
architectures. The theoretical foundation for these approaches is the
Johnson-Lindenstrauss (JL) lemma \cite{johnson1984}, which guarantees that
pairwise distances can be approximately preserved after projection into a
lower-dimensional space with high probability.

Despite their efficiency, conventional dense Gaussian random projections
still impose non-trivial computational overhead, and recent analyses have
shown that random projections can degrade node similarity measures under
certain graph structures \cite{tadic2024}. Sparse and structured projection
matrices, as established by Achlioptas \cite{achlioptas2003}, can reduce
this cost while maintaining JL guarantees.

Building on these observations, this work addresses the following research
question: \emph{Can hybrid structured random projections---combining sparse
structural embeddings with node feature information---outperform a dense
Gaussian baseline in terms of downstream node classification and link
prediction accuracy, while remaining computationally tractable?}

The main contributions of this paper are:
\begin{itemize}
    \item A clean, modular implementation of three projection strategies
          (Gaussian, Striped, Hybrid) within a unified FastRP-style pipeline.
    \item A systematic evaluation across four node-classification datasets
          and three link-prediction datasets, with Optuna-based
          hyperparameter tuning for fair comparison.
    \item Empirical evidence that the hybrid structure-feature projection
          substantially and consistently outperforms the Gaussian baseline on
          feature-rich graphs.
    \item An analysis of the conditions under which structured sparse
          projections succeed and where they remain sensitive.
\end{itemize}

The remainder of the paper is organized as follows. Section~\ref{sec:related}
reviews related work. Section~\ref{sec:method} describes the methodology.
Section~\ref{sec:setup} covers the experimental setup. Section~\ref{sec:results}
presents and discusses the results. Section~\ref{sec:conclusion} concludes.

% =========================================================================
\section{Related Work}
\label{sec:related}

\subsection{Scalable Graph Embedding}

Projection-free embedding methods that rely on random walks---including
DeepWalk \cite{perozzi2014} and Node2vec \cite{grover2016node2vec}---apply
Skip-Gram language modelling to graph traversals to learn latent node
representations. DeepWalk achieves up to 10\% higher F1 than spectral
baselines on sparse-label benchmarks, but the walk generation step is
expensive and difficult to parallelise at scale. Node2vec extends this with
biased walks controlled by return and in-out parameters, yielding greater
flexibility in capturing breadth-first versus depth-first neighbourhood
structure \cite{grover2016node2vec}.

To avoid walk generation, matrix-factorisation approaches such as
LINE \cite{tang2015line} optimise explicit proximity objectives for first-
and second-order proximity, while NetSMF \cite{qiu2019netsmf} approximates
the implicit matrix objective of DeepWalk through sparse matrix factorisation,
demonstrating that sparse operations are essential for billion-scale graphs.
HOPE \cite{ou2016atp} addresses directed graphs by preserving asymmetric
transitivity via high-order proximity measures such as the Katz index and
rooted PageRank, showing that orientation-aware proximity cannot be
recovered by symmetric-proximity methods alone.

Projection-based approaches eliminate both walk generation and matrix
decomposition. RandNE \cite{zhang2018randne} uses iterative random
projections to preserve high-order proximity without expensive
factorisations, scaling to billion-edge graphs. SketchNE
\cite{xie2023sketchne} combines sketching with randomised SVD, embedding
billion-node networks in under one hour and outperforming walk-based methods
in both speed and quality. FastRP \cite{chen2019fast} reduces density
further with very sparse projections and iterative propagation, achieving
Macro-F1 scores competitive with walk-based methods while reducing runtime
by orders of magnitude, confirming that projection-based pipelines can
match walk-based quality at a fraction of the cost.

\subsection{Structured, Sparse, and Hybrid Random Projections}

The theoretical foundation for sparse projection matrices was established
by Achlioptas \cite{achlioptas2003}, who proved that entries drawn from
$\{-1, 0, +1\}$ with up to two-thirds zero density preserve pairwise
distances under the Johnson-Lindenstrauss lemma \cite{johnson1984}.
This result implies that sparsity does not fundamentally degrade embedding
quality and motivates the construction of structured sparse matrices that
are both fast to apply and memory-efficient.

Yahaya et al. \cite{yahaya2025hrp} formalised Hybrid Random Projections
(HRP) as a mixture $\mathbf{H} = \alpha\mathbf{R} + \beta\mathbf{S}$,
where $\mathbf{R}$ is a classical Gaussian matrix, $\mathbf{S}$ encodes
structural constraints, and the scalars $\alpha, \beta$ govern the
randomness--structure balance. Across several high-dimensional benchmark
datasets, this formulation consistently reduces geometric distortion relative
to purely Gaussian projections. Hierarchical feature selection combined
with random projection \cite{wang2019hierarchical} further refines
projection pipelines by identifying salient dimensions before projection,
reducing noise and memory overhead without expensive matrix inversion.

In the GNN context, RpHGNN \cite{hu2021rphgnn} compresses heterogeneous
graph tensors via random projection for efficient mini-batch training,
reporting speedups exceeding 230\% over conventional GNN training.
RAP-GNN \cite{bui2023rap} replaces all learnable message-passing weights
with diagonal random matrices sampled on-the-fly, cutting training time by
a factor of six and memory by a factor of three while achieving accuracy
comparable to fully trained GNNs---challenging the assumption that explicit
parameter learning is always required. Random projection forest
initialisation \cite{alshammari2020rpforest} improves GCN embedding quality
by weighting edges according to co-occurrence frequency across multiple
projection trees, producing higher-quality neighbourhood graphs than
standard $k$-NN constructions.

Despite these advances, Tad\'{i}\'{c} et al. \cite{tadic2024} document
pathological cases in which random projections degrade node similarity
and ranking quality, particularly for nodes at the extremes of the degree
distribution, underscoring that projection design must be sensitive to
graph topology.

\subsection{Research Gap}

Existing work on projection-based graph embedding covers a range of
projection designs: FastRP \cite{chen2019fast} uses very sparse
projections with $O(\sqrt{m})$ sparsification, RandNE \cite{zhang2018randne}
uses Gaussian random projection, and SketchNE \cite{xie2023sketchne}
combines randomised SVD with sketching. None of these, nor methods that
integrate random projections into GNN training loops (RpHGNN, RAP-GNN),
have conducted a controlled, multi-dataset comparison of
(a) a standard dense Gaussian baseline,
(b) a structured sparse projection inspired by Achlioptas \cite{achlioptas2003},
and (c) a hybrid projection that combines structural embedding with node
feature information, all within the same iterative propagation pipeline
and with identical hyperparameter tuning infrastructure.
The present work fills this gap and provides the first systematic evidence
of the conditions under which each projection type is beneficial.

% =========================================================================
\section{Methodology}
\label{sec:method}

\subsection{Problem Formulation}

Let $G = (V, E)$ be an undirected graph with $N = |V|$ nodes and adjacency
matrix $\mathbf{A} \in \{0,1\}^{N \times N}$. The row-normalized transition
matrix is $\mathbf{T} = \mathbf{D}^{-1}\mathbf{A}$, where $\mathbf{D}$ is
the diagonal degree matrix. Some datasets also provide a node feature matrix
$\mathbf{X} \in \mathbb{R}^{N \times F}$. The goal is to learn an embedding
matrix $\mathbf{Z} \in \mathbb{R}^{N \times d}$ ($d \ll N$) suitable for
downstream node classification or link prediction.

\subsection{Base Method: FastRP}
\label{subsec:fastrp-base}

Our pipeline builds directly on FastRP \cite{chen2019fast}, which establishes
the two-component design paradigm for scalable graph embedding: (1)
constructing a node similarity matrix that preserves high-order graph
proximity, and (2) applying an optimization-free dimension reduction method
to that matrix.

\paragraph{Similarity matrix.}
FastRP \cite{chen2019fast} observes that the $k$-step transition matrix
$\mathbf{T}^k$ converges entrywise towards $d_j / 2m$ as $k$ grows---a
consequence of the Perron--Frobenius theorem applied to the symmetrized
transition matrix \cite{chen2019fast}. Because many real-world graphs are
scale-free, this produces a heavy-tailed entry distribution that harms
dimension reduction \cite{chen2019fast}. To correct for this, FastRP
\cite{chen2019fast} applies a degree-based normalization derived from a
Tukey power transform \cite{tukey1957}: each column $j$ of $\mathbf{T}^k$
is scaled by $(d_j / 2m)^{\beta}$ for a tunable exponent $\beta \in [-1, 0]$,
yielding the corrected matrix
\begin{equation}
    \widetilde{\mathbf{T}}^k_{ij}
        = \mathbf{T}^k_{ij} \left(\frac{d_j}{2m}\right)^{\beta}.
    \label{eq:fastrp-norm}
\end{equation}
This normalization downweights the influence of high-degree hub nodes and
can be expressed as right-multiplication by a diagonal matrix $\mathbf{L}$,
enabling fully iterative computation \cite{chen2019fast}.

\paragraph{Very sparse random projection.}
For dimension reduction, FastRP \cite{chen2019fast} draws the projection
matrix $\mathbf{R} \in \mathbb{R}^{N \times d}$ from the very sparse
distribution of Li et al. \cite{li2006very}, where each
entry is independently drawn from
\begin{equation}
    R_{ij} = \begin{cases}
        +\!\sqrt{s}, & \text{with probability } \tfrac{1}{2s},\\
        0,           & \text{with probability } 1 - \tfrac{1}{s},\\
        -\!\sqrt{s}, & \text{with probability } \tfrac{1}{2s},
    \end{cases}
    \label{eq:very-sparse}
\end{equation}
with $s = \sqrt{m}$ recommended, achieving $\sqrt{m}$-fold speedup over
dense Gaussian projection while preserving Johnson--Lindenstrauss
guarantees \cite{dasgupta1999,li2006very}.

\paragraph{Iterative computation and weighted combination.}
A key algorithmic contribution of FastRP \cite{chen2019fast} is the
exploitation of matrix associativity to avoid forming $\mathbf{T}^k$
explicitly. Starting from $\mathbf{N}_1 = \mathbf{T}\mathbf{L}\mathbf{R}$,
each subsequent order is obtained as $\mathbf{N}_i = \mathbf{T}\mathbf{N}_{i-1}$
at cost $O(md)$ per step. The final embedding is a weighted sum over
propagation orders:
\begin{equation}
    \mathbf{N} = \sum_{i=1}^{k} \alpha_i \mathbf{N}_i,
    \label{eq:fastrp-combine}
\end{equation}
where $\alpha_1, \ldots, \alpha_k$ are tunable scalar weights. This
iterative scheme gives FastRP \cite{chen2019fast} overall time complexity
$O((N + m) k d)$, linear in the graph size.

\subsection{Our Implementation and Adaptations}
\label{subsec:implementation}

We implement the above FastRP pipeline with the following
adaptation to the normalization step. Rather than applying the exact
$(d_j/2m)^{\beta}$ column scaling from Equation~\eqref{eq:fastrp-norm},
we scale each row of the initial representation by $d_i^{\alpha}$ (fixed
at $\alpha = -0.6$ in all experiments). FastRP's normalization
\eqref{eq:fastrp-norm} scales column $j$ by $(d_j/2m)^\beta$, which
reduces the contribution that high-degree \emph{target} nodes make to
all embeddings. Our row scaling reduces the \emph{source} node $i$'s
contribution to its own initial representation; both choices discourage
high-degree nodes from dominating the subsequent aggregation, but they
operate in different directions. The exponent $\alpha$ serves the same
regularizing purpose as FastRP's $\beta$ but is kept fixed (rather than
tuned) to isolate the effect of the projection strategy.

Additionally, we apply $\ell_2$-normalization to each intermediate
representation $\mathbf{N}^{(t)}$ before accumulation, and apply a final
per-dimension $z$-score standardization to the aggregated embedding
$\mathbf{Z}$. Neither of these post-processing steps appears in the
original FastRP paper \cite{chen2019fast}; they are introduced to improve
numerical stability and conditioning of the downstream logistic regression
classifier.

Formally, the pipeline as implemented proceeds as follows:

\paragraph{Step 1: Initial projection.}
Construct the initial representation $\mathbf{N}^{(0)} \in \mathbb{R}^{N \times d}$
according to the chosen projection strategy (Section~\ref{subsec:strategies}).

\paragraph{Step 2: Degree-based row scaling (adapted from FastRP \cite{chen2019fast}).}
Scale each row $i$ of $\mathbf{N}^{(0)}$ by $d_i^{\alpha}$, where $\alpha = -0.6$.
This reduces hub-node dominance, analogously to FastRP's column normalization
\eqref{eq:fastrp-norm}.

\paragraph{Step 3: Iterative neighborhood propagation (inherited from FastRP \cite{chen2019fast}).}
For $t = 1, \ldots, k$:
\begin{equation}
    \mathbf{N}^{(t)} = \mathbf{T}\, \mathbf{N}^{(t-1)}.
    \label{eq:propagation}
\end{equation}
The embedding is accumulated as a weighted sum of $\ell_2$-normalized
intermediate representations:
\begin{equation}
    \mathbf{Z}_\text{raw} = \sum_{t=1}^{k} w_t \cdot
        \frac{\mathbf{N}^{(t)}}{\|\mathbf{N}^{(t)}\|_2}.
    \label{eq:weighted-sum}
\end{equation}

\paragraph{Step 4: Final standardization (our addition).}
The aggregated embedding is $z$-score standardized per dimension to produce
the final output $\mathbf{Z}$.

\subsection{Novel Projection Strategies}
\label{subsec:strategies}

The central contribution of this work is the replacement of FastRP's
very sparse projection \cite{li2006very,chen2019fast} with three
alternative strategies for constructing $\mathbf{N}^{(0)}$, enabling a
controlled study of how projection design affects downstream performance.

\subsubsection{Gaussian Baseline}
The projection matrix entries are drawn i.i.d. from a standard normal
distribution:
\begin{equation}
    R_{ij} \sim \mathcal{N}(0, 1).
\end{equation}
Dense Gaussian projection is the classical reference for JL-type
embeddings \cite{dasgupta1999} and serves as the baseline against which
sparser and hybrid strategies are compared. It is denser than FastRP's
very sparse projection \cite{chen2019fast,li2006very} and incurs higher
memory and compute cost.

\subsubsection{Striped Structured Sparse Projection (Our Contribution)}
We propose a structured sparse projection in which the $d$ output
dimensions are partitioned into non-overlapping groups of size $g$
(a tunable hyperparameter). For each node $i$ and each group, exactly one
column within the group receives a non-zero entry drawn uniformly from
$\{-1, +1\}$, scaled by $\sqrt{g}$ to preserve the variance of the
Gaussian baseline:
\begin{equation}
    R_{ij} \in \{-\sqrt{g},\; 0,\; +\sqrt{g}\},
    \quad \text{density} = \frac{1}{g}.
    \label{eq:striped}
\end{equation}
The group constraint enforces a regular non-zero pattern (``stripes'')
across the projection matrix, in contrast to the independent Bernoulli
draws of Achlioptas \cite{achlioptas2003} or Li et al.\ \cite{li2006very}.
Larger $g$ increases sparsity. JL guarantees are preserved at any density
above $\Omega(\log N / d)$ \cite{achlioptas2003,johnson1984}. Relative
to FastRP's very sparse projection, which achieves $s = \sqrt{m}$-fold
speedup through fully independent entries, the striped design trades
independence across entries within a group for a deterministic coverage
pattern that ensures each group of $g$ dimensions always contains exactly
one non-zero entry per node---a property we introduce to prevent any
dimension block from being entirely zeroed out.

\subsubsection{Hybrid Structure-Feature Projection (Our Contribution)}
When node features $\mathbf{X}$ are available, we partition the $d$
embedding dimensions into a structural component ($d_\text{struct} =
\lfloor 3d/4 \rfloor$) filled by the striped projection, and a feature
component ($d_\text{feat} = d - d_\text{struct}$) derived from the node
features through a random Gaussian compression:
\begin{equation}
    \mathbf{R}_\text{feat} = \mathbf{X}\, \mathbf{P},
    \quad \mathbf{P} \in \mathbb{R}^{F \times d_\text{feat}},
    \quad P_{ij} \sim \mathcal{N}\!\left(0,\, \tfrac{1}{F}\right).
    \label{eq:hybrid-feat}
\end{equation}
The two components are concatenated to form the initial representation:
\begin{equation}
    \mathbf{N}^{(0)} = \bigl[\mathbf{R}_\text{struct} \;\|\; \mathbf{R}_\text{feat}\bigr].
    \label{eq:hybrid-concat}
\end{equation}
This design extends the FastRP pipeline \cite{chen2019fast}---which is
purely topology-driven---by fusing node attribute information into the
initial projection. The subsequent propagation steps then propagate both
the structural and semantic signals jointly across the graph. The 3:1
structure-to-feature dimension ratio was chosen empirically; ablating
this ratio is left as future work.

\subsection{Downstream Tasks}

\paragraph{Node Classification.}
Embeddings are passed to a one-vs-rest logistic regression classifier
(\texttt{liblinear} solver) trained on 80\% of nodes and evaluated on the
remaining 20\%, measuring Macro-F1 score. For ogbn-arxiv, standard
accuracy on the OGB split is also reported.

\paragraph{Link Prediction.}
Edges are split into training (85\%), validation (5\%), and test (10\%)
sets. For each positive test edge, 500 negative edges are sampled. Link
scores are computed as the dot product of node embeddings, and performance
is measured by AUC and MRR@10 averaged over five independent random seeds.
The model is tuned on the validation set and then evaluated on the test set.

% =========================================================================
\section{Experimental Setup}
\label{sec:setup}

\subsection{Datasets}

Table~\ref{tab:datasets} summarizes the benchmark datasets used.

\begin{table}[ht]
\centering
\caption{Benchmark Datasets}
\label{tab:datasets}
\begin{tabular}{lrrrl}
\toprule
\textbf{Dataset} & \textbf{Nodes} & \textbf{Edges} & \textbf{Features} & \textbf{Task} \\
\midrule
BlogCatalog & 10,312 & 333,983 & None & NC \\
Flickr      & 89,250 & 899,756 & None & NC \\
Cora        & 2,708  & 5,429   & 1,433 & NC / LP \\
ogbn-arxiv  & 169,343 & 1,166,243 & 128 & NC \\
Citeseer    & 3,327  & 4,732   & 3,703 & LP \\
PubMed      & 19,717 & 44,338  & 500  & LP \\
\bottomrule
\end{tabular}
\end{table}

\textit{NC = Node Classification; LP = Link Prediction.}

BlogCatalog and Flickr are social/content-sharing networks loaded from
standard \texttt{.mat} files. Cora, Citeseer, and PubMed are citation
network benchmarks downloaded via PyTorch Geometric. ogbn-arxiv is loaded
from the Open Graph Benchmark \cite{hu2020ogb}.

\subsection{Implementation}

The framework is implemented in Python using PyTorch for tensor operations
and scipy for sparse matrix handling. The \texttt{FastRP} class supports
all three projection types via a single \texttt{projection\_type} argument.
All experiments are executed with PyTorch \texttt{no\_grad} (no
backpropagation), so the approach is strictly unsupervised at the embedding
stage.

\subsection{Hyperparameter Tuning}

Hyperparameters are tuned using Optuna \cite{optuna} (TPE sampler) with 50
trials per experiment, optimizing Macro-F1 (node classification) or MRR@10
(link prediction) on a held-out validation split. The tuning space is:

\begin{itemize}
    \item \textbf{Embedding dimension} $d \in \{64, 128, 256, 512\}$
    \item \textbf{Window size} $k \in [1, 5]$
    \item \textbf{$\ell_2$ normalization} $\in \{\text{True}, \text{False}\}$
    \item \textbf{Group size} $g \in [2, 10]$ (striped and hybrid only)
\end{itemize}

The following parameters are fixed across all runs: $\alpha = -0.6$,
propagation weights $\mathbf{w} = [1.0, 1.0, 7.81, 45.28]$ (derived from
the BlogCatalog baseline), and \texttt{input\_matrix = `trans'} (row-normalized
transition matrix). Final test-set results for link prediction are averaged
over five random seeds (42, 100, 200, 300, 400) using the best
hyperparameters from the validation phase.

% =========================================================================
\section{Results and Discussion}
\label{sec:results}

\subsection{Node Classification}

\subsubsection{Cora}

Table~\ref{tab:nc_cora} reports Macro-F1 scores on Cora. Both structured
projection variants outperform the Gaussian baseline. The hybrid model
achieves the highest score, demonstrating that incorporating node feature
information (1,433-dimensional bag-of-words vectors) into the initial
projection leads to measurably richer embeddings on this citation network.

\begin{table}[ht]
\centering
\caption{Node Classification --- Cora (Macro-F1)}
\label{tab:nc_cora}
\begin{tabular}{lc}
\toprule
\textbf{Method} & \textbf{Macro-F1} \\
\midrule
Gaussian Baseline         & 0.7232 \\
Striped Structured        & 0.7356 \\
\textbf{Hybrid}           & \textbf{0.7541} \\
\bottomrule
\end{tabular}
\end{table}

\subsubsection{ogbn-arxiv}

Table~\ref{tab:nc_arxiv} reports train, validation, and test accuracy on
ogbn-arxiv. The Gaussian and striped projections achieve comparable scores
($\sim$0.52 test accuracy), performing poorly on this large feature-rich
graph. The hybrid projection achieves a test accuracy of 0.6891---a gain
of 16.8 percentage points over each of the structure-only baselines---
highlighting that the 128-dimensional paper embedding features in ogbn-arxiv
carry strong discriminative signal that the structure-only methods cannot
exploit.

\begin{table}[ht]
\centering
\caption{Node Classification --- ogbn-arxiv (Accuracy)}
\label{tab:nc_arxiv}
\begin{tabular}{lccc}
\toprule
\textbf{Method} & \textbf{Train} & \textbf{Val} & \textbf{Test} \\
\midrule
Gaussian Baseline   & 0.5149 & 0.5080 & 0.5211 \\
Striped HRP         & 0.5115 & 0.5063 & 0.5184 \\
\textbf{Hybrid HRP} & \textbf{0.7468} & \textbf{0.6970} & \textbf{0.6891} \\
\bottomrule
\end{tabular}
\end{table}

\subsubsection{Flickr}

Table~\ref{tab:nc_flickr} reports Macro-F1 on Flickr, which has no node
features, so only Gaussian and striped projections are directly comparable.
The striped projection (0.1961) slightly underperforms the Gaussian baseline
(0.2030). This result indicates that striped structured sparsity is not
universally beneficial and is sensitive to the dataset's topology and label
distribution. Flickr is a large social graph (89K nodes) with multi-label
classification; the relatively low absolute Macro-F1 values reflect the
difficulty of the task.

\begin{table}[ht]
\centering
\caption{Node Classification --- Flickr (Macro-F1)}
\label{tab:nc_flickr}
\begin{tabular}{lc}
\toprule
\textbf{Method} & \textbf{Macro-F1} \\
\midrule
Gaussian Baseline   & 0.2030 \\
Striped Structured  & 0.1961 \\
Hybrid              & \textit{N/A (no node features)} \\
\bottomrule
\end{tabular}
\end{table}

\subsection{Link Prediction}

Table~\ref{tab:lp_results} reports link prediction performance (AUC and
MRR@10) on Cora, Citeseer, and PubMed using the best hyperparameters
selected by Optuna on the validation split. Results for the Gaussian baseline
and the Hybrid projection are reported; the striped-only variant was not
evaluated separately in the link prediction experiments.

\begin{table}[ht]
\centering
\caption{Link Prediction Performance (AUC / MRR@10)}
\label{tab:lp_results}
\begin{tabular}{llll}
\toprule
\textbf{Variant} & \textbf{Cora} & \textbf{Citeseer} & \textbf{PubMed} \\
\midrule
Gaussian  & 84.50\% / 43.29\% & 77.19\% / 39.08\% & 78.22\% / 43.49\% \\
\textbf{Hybrid}    & \textbf{87.90\%} / \textbf{44.08\%} & \textbf{81.08\%} / \textbf{42.17\%} & \textbf{85.09\%} / 15.24\% \\
\bottomrule
\end{tabular}
\end{table}

Table~\ref{tab:lp_time} reports mean embedding CPU time per variant. Values
are wall-clock process times measured by \texttt{time.process\_time()}.

\begin{table}[ht]
\centering
\caption{Embedding CPU Time for Link Prediction (seconds)}
\label{tab:lp_time}
\begin{tabular}{lrrr}
\toprule
\textbf{Variant} & \textbf{Cora} & \textbf{Citeseer} & \textbf{PubMed} \\
\midrule
Gaussian & 0.090 & 0.057 & 1.434 \\
Hybrid   & 0.122 & 0.090 & 0.970 \\
\bottomrule
\end{tabular}
\end{table}

\subsection{Discussion}

\subsubsection{Why Hybrid Projections Help on Feature-Rich Graphs}

The most striking finding of this study is the 16.8 percentage-point
improvement in test accuracy achieved by the hybrid projection on ogbn-arxiv
over both structure-only baselines (Table~\ref{tab:nc_arxiv}). This result
deserves careful interpretation. The Gaussian and striped projections assign
random coordinates to each node before graph propagation; because these
coordinates carry no semantic content, the only signal available to
downstream classification comes from structural proximity after aggregation.
Graph topology alone is insufficient to distinguish the 40 subject-area
categories in ogbn-arxiv, where many papers from different fields are
cited by similar hub papers.

The hybrid projection allocates $d_\text{feat} = d/4$ of the embedding budget
to a random projection of the 128-dimensional node features, which are
paper title and abstract embeddings generated by skip-gram word vectors.
Even after a single feature projection step, these dimensions encode
semantic similarity: two papers in the same sub-field will have nearby
feature projections, and the subsequent graph propagation reinforces this
alignment by aggregating projections from structurally similar neighbours.
This combination of semantic initialisation and structural smoothing produces
a considerably more discriminative embedding than either component alone, a
mechanism that echoes the motivation for combining explicit node attributes
with neighbourhood information in methods such as GraphSAGE \cite{khemani2024}
and modern GNN surveys \cite{gkarmpounis2024}.

A similar, if smaller, benefit holds on Cora (Macro-F1: 0.7541
vs. 0.7232 Gaussian), where bag-of-words features provide lexically
distinct signals for the seven paper categories. The fact that the gain
on Cora is smaller than on ogbn-arxiv is consistent with this interpretation:
Cora's graph is small and tightly clustered, so structural propagation
already recovers much of the class signal even without feature information.

\subsubsection{Why Structured Sparsity Is Dataset-Sensitive}

The striped projection outperforms the Gaussian baseline on Cora (0.7356
vs. 0.7232) but trails it on Flickr (0.1961 vs. 0.2030). This divergence
sheds light on the conditions under which structured sparsity is beneficial.

On Cora, the graph is sparse and exhibits clear community structure aligned
with paper categories. In this setting, the group-wise constraint of the
striped projection---exactly one non-zero entry per group of $g$ dimensions---
acts as a form of regularisation that reduces inter-dimension correlation in
the initial representation, potentially improving the geometry of the
embedding before propagation. The JL lemma guarantees that such sparse
matrices preserve pairwise distances \cite{johnson1984,achlioptas2003}, and
the structured sparsity may reduce the effective noise level in the initial
embedding, giving graph propagation a cleaner signal to amplify.

On Flickr, the graph is much denser (approximately 10$\times$ more edges
per node than Cora) and the task involves 195-class multi-label
classification, where the decision boundaries are spread diffusely across
the feature space. In such settings, the group-wise constraint of the
striped matrix may introduce correlated dead zones across embedding
dimensions---if a group's non-zero entry happens to be uninformative for a
particular node, all $g-1$ remaining dimensions carry zero signal for that
node. This pathology is analogous to the degradation documented by
Tad\'{i}\'{c} et al. \cite{tadic2024} for degree-skewed graphs: structure
imposed on a projection that was not designed for the graph's topology can
actively harm embedding quality. Together these observations suggest that
the group size $g$ should be calibrated to the graph's density and the
task's label complexity.

\subsubsection{Link Prediction: Hybrid Consistently Outperforms Gaussian}

Across all three citation benchmarks (Cora, Citeseer, PubMed), the hybrid
projection achieves higher AUC than the Gaussian baseline
(Table~\ref{tab:lp_results}), with improvements of 3.40, 3.89, and
6.87 percentage points respectively. The node features in these datasets
(bag-of-words for Cora and Citeseer; TF-IDF for PubMed) provide
semantically meaningful similarity information: two papers that share
many keywords are likely to be related and, therefore, potentially linked.
By incorporating feature similarity into the initial projection, the hybrid
embedding makes the dot-product link score a better approximation of true
citation probability than a structure-only score.

The one anomalous result is PubMed's MRR@10 for the hybrid (15.24\%
vs. 43.49\% for Gaussian). AUC and MRR@10 measure complementary aspects of
ranking quality: AUC measures global discrimination while MRR@10 rewards
having at least one true positive in the top-10 ranked candidates per query.
The large MRR@10 gap suggests that, despite achieving high overall AUC, the
hybrid embedding fails to place true positive links in the very top positions
for many PubMed query nodes. We hypothesize that because PubMed's feature
space (TF-IDF over biomedical vocabulary) is high-dimensional and sparse,
the feature projection---which uses a Gaussian matrix scaled by $1/F$---may
require more than $d/4$ dimensions to faithfully represent biomedical
relatedness, leading to imprecise top-$k$ ranking even when global AUC is
better. While this interpretation motivates future experiments to vary the
structure-to-feature split ratio, further investigation is required to
isolate the exact cause of this anomaly.

\subsubsection{Computational Efficiency}

A key design goal of projection-based embeddings is to eliminate the
backpropagation cost of end-to-end GNNs. All methods in this study
satisfy this goal: no gradients are computed, and embedding generation
requires a single matrix multiplication per propagation step. The embedding
CPU times reported in Table~\ref{tab:lp_time} confirm that the evaluated
Gaussian and hybrid variants finish in well under two seconds even on
PubMed (19.7K nodes).

A notable result is that the hybrid projection is \emph{faster} than the
Gaussian baseline on PubMed (0.970\,s vs. 1.434\,s). This counterintuitive
outcome arises because the hybrid replaces $d/4$ of the full-graph structural
projection dimensions with a feature projection step; the feature projection
is a dense matrix multiply of size $N \times F$ times $F \times d/4$, but
$F = 500$ is small relative to $N = 19{,}717$, making this operation
cheaper than the corresponding graph propagation over 500 dimensions. On
smaller graphs (Cora: 2,708 nodes, Citeseer: 3,327 nodes), this substitution
is less impactful and the hybrid incurs a modest overhead of at most 35\,ms
due to the additional feature projection multiply. These results are
consistent with the memory-efficiency findings reported for RAP-GNN
\cite{bui2023rap} and RpHGNN \cite{hu2021rphgnn}: carefully designed
projection operators can simultaneously improve accuracy \emph{and} reduce
runtime relative to naive dense baselines.

A full comparison with end-to-end GNN runtimes (GCN, GraphSAGE, GAT) is
beyond the scope of this study---such models require multiple training
epochs and GPU infrastructure---but the single-forward-pass nature of the
present approach suggests it could serve as a fast initialisation stage
before fine-tuning, following the spirit of RP-forest GCN initialisation
\cite{alshammari2020rpforest}.

\subsubsection{Limitations and Threats to Validity}

\textbf{Absence of strong GNN baselines.}
This study does not compare against end-to-end trained GNNs such as GCN,
GraphSAGE \cite{khemani2024}, or GAT. The primary claim is a controlled
comparison \emph{across projection types}, not a claim that projection-based
methods are superior to learned methods in absolute terms. Including GNN
baselines is important future work for situating these results in the
broader landscape.

\textbf{Fixed structure-to-feature split.}
The 3:1 structure-to-feature dimension ratio in the hybrid projection was
chosen empirically without systematic ablation. Different ratios may be
optimal for different datasets or tasks, and the anomalous PubMed MRR@10
result specifically suggests that a larger feature share may be beneficial
for high-dimensional sparse feature spaces.

\textbf{Single projection trial.}
Unlike the link prediction pipeline (which averages over five random seeds),
node classification results for a given hyperparameter configuration reflect
a single random seed for the projection matrix. Variability across projection
seeds could affect Macro-F1 by a small amount, particularly for Flickr where
the margin between Gaussian and striped is only 0.0069.

\textbf{Homogeneous graph assumption.}
All datasets used are homogeneous (single node and edge type). The hybrid
projection has not been evaluated on heterogeneous graphs, where the
distinction between node types would require type-conditioned projection
strategies as used in RpHGNN \cite{hu2021rphgnn}.

\textbf{Transductive setting.}
All node classification experiments operate in the transductive setting
(all nodes visible at embedding time). Inductive extensions, where test
nodes are unseen during embedding, require additional study.

% =========================================================================
\section{Conclusion}
\label{sec:conclusion}

This paper investigated the effect of projection matrix design on the
quality of graph embeddings within a FastRP-style framework. Three
strategies were compared for node classification---a dense Gaussian baseline,
a striped structured sparse projection, and a hybrid structure-feature
projection---and two strategies (Gaussian and hybrid) were compared for
link prediction, evaluated across seven benchmark datasets.

The results establish that the hybrid structure-feature projection
consistently achieves the best predictive performance on datasets with
informative node attributes, with particularly large gains on the
large-scale ogbn-arxiv benchmark (16.8 percentage point improvement in
test accuracy). The striped structured projection demonstrates competitive
performance on citation graphs, supporting the viability of sparse
structured projections as a scalable alternative to dense Gaussian
projections. However, its sensitivity on the Flickr social graph
underscores the importance of dataset-aware hyperparameter tuning,
particularly the group size $g$.

Future work should explore: (i) a systematic ablation of the
structure-to-feature dimension ratio in the hybrid projection; (ii)
integration of these projection strategies into full GNN architectures as
a preprocessing or regularization component; (iii) extension to
heterogeneous and temporal graphs; and (iv) a detailed runtime and memory
efficiency characterization across varying graph sizes.

% =========================================================================
\section*{Author Contributions}

H.A.D.T.N. Haththella contributed to the design of the projection
strategies, implementation of the striped structured projection, and the
node classification experiments.
L.R.S.R. Premasiri contributed to the implementation of the hybrid
projection, the link prediction pipeline, and hyperparameter tuning with
Optuna.
A.N.T. Somathilaka contributed to dataset preparation, data loading
infrastructure, and experimental analysis.
All authors participated in drafting and revising the manuscript.
Dr. Sampath Deegalla supervised the project.

% NOTE TO AUTHORS: Please confirm and update the above contribution
% statements to accurately reflect individual roles.

% =========================================================================
\bibliographystyle{IEEEtran}
\bibliography{ref}

\end{document}
